\section{Threats to Validity}\label{sec:threats}
\paragraph*{Construct} Our consistency metrics are the system's own checks (schema placement, owner scope, foreign references, registry lockstep) plus an authorship marker for cross-project writes; violations outside these classes are not counted. Task success for the LLM runs is a module plus a passing test that \emph{we} execute, not the agents' claims.
\paragraph*{Internal} The scripted workload (RQ1) replays the error classes observed in the system's history with fixed proportions. It demonstrates the mechanism, not the rate at which real agents err; RQ2 measures that. The before-and-after comparison changes one thing at a time only in part: P2a and P2b were measured together after the watch stage, but the hook blocked nothing in the normal workload, so the placement gain is attributable to P2a. The adversarial variant is a single run per condition at $N=5$. The hook's shell parser is best-effort, and writes from inside a script are left to the audit. All runs used the same machine while the owner's production instance and scheduled security agents were active; we record their load and report ratios rather than absolute times. Cloud and remote tools were blocked by design, so deployment-to-VM steps of the owner's tests were run locally (Section~\ref{sec:method}). The registry database is a private copy restored from a read-only dump. Two weaknesses of the studied version lie outside our consistency metrics: per-project deltas were generated but not always saved and registered, and the truth gate did not yet cover the reviewer's verdict, so a run could report success while its review asked for a fix. Neither affects the reported placement, violation or task-success figures, which are taken from disk, from the system's own checks and from tests that we run ourselves.
\paragraph*{External} This is one system, built mainly by one engineer with agents, on one laptop. Results may not transfer to other organisations; the mechanisms (one door, registry-only database, audit) are, however, framework-independent.
\paragraph*{Reliability} We repeated every run: 5 seeds for RQ1; 3 repetitions per $N$ before and after for RQ2; 3 repetitions for the owner's scripted tests and for RQ4. The LLM-driven owner tests ran once per state. One invalid run is excluded and reported: our own regression run reset the database under it. A replication package with all scripts, raw results and the decision log is provided.
